\section{Volume-normalized Berger links}
\label{sec:links}

We need angular metrics that retain fixed finite-dimensional harmonic spaces
while varying their spectra. This section derives their curvature and angular
operators; the next section selects the directions that will be cycled.

Fix an odd integer \(m=2s-1\), with \(s\geq4\), and let \(g_0\) be the
unit round metric on \(S^m\subset\R^{2s}\).
An orthogonal complex structure is a real linear map \(J\) satisfying
\(J^2=-\Id\) and \(J^{\mathsf T}J=\Id\).
It determines the unit round Killing field and its dual one-form
\[
 \xi_J(x)=Jx,\qquad \eta_J=g_0(\xi_J,\cdot).
\]
Write
\[
 P_J=\xi_J\otimes\eta_J,\qquad D_Ju=du(\xi_J).
\]
Thus \(P_J\) is the projection onto the Hopf direction.
For \(a,q>0\), define
\[
 \widetilde g_{q,J}=g_0+(q-1)\eta_J^2,\qquad
 G(a,q,J)=a^2q^{-1/m}\widetilde g_{q,J}.
\]
The factor \(q^{-1/m}\) makes the volume independent of \(q\).
All orthogonal complex structures will be allowed; we do not restrict
them to one component of their parameter space.

This volume normalization and the corresponding Laplacian decomposition
are classical; see Tanno~\cite[Sections~2--4]{Tanno}.
His parameter \(t\) corresponds to \(q=t^m\) when \(a=1\).
Classical curvature calculations for these spheres appear in
Bourguignon--Karcher~\cite[Section~6.6]{BourguignonKarcher}.
We record the formulas and their direct proofs in our parameters below.

For \(l\geq0\), let \(V_l\) be the real vector space of restrictions to
\(S^m\) of homogeneous harmonic polynomials of degree \(l\) on
\(\R^{m+1}\).  Equip \(V_l\) with the round \(L^2\) inner product and put
\[
 N_l=\dim_{\R}V_l
     =\binom{m+l}{m}-\binom{m+l-2}{m},
 \qquad B_l=l(l+m-1),
\]
where the second binomial coefficient is zero for \(l<2\).
For completeness, on homogeneous polynomials give distinct monomials
orthogonal inner products with
\(\langle x^\alpha,x^\alpha\rangle=\alpha!\).
The adjoint of the Euclidean polynomial Laplacian from degree \(l\) to
degree \(l-2\) is multiplication by \(|x|^2\), which is injective.
The Laplacian is therefore surjective, giving the dimension formula.
Restriction to the sphere is injective on homogeneous polynomials.
Euclidean polar coordinates give
\[
 \Delta_{g_0}^{+}\big|_{V_l}=B_l\Id,
 \qquad \Delta_h^{+}:=-\operatorname{div}_h\nabla_h.
\]
Rotations preserve \(V_l\).  In particular, \(D_J\) preserves \(V_l\)
and is skew-adjoint for its round inner product.

\begin{lemma}\label{lem:berger-metrics}
The volume form of \(G(a,q,J)\) is \(a^m\,d\operatorname{vol}_{g_0}\).
Its Ricci endomorphism has eigenvalues
\[
 a^{-2}q^{1/m}\bigl((m-1)-2(q-1)\bigr)
 \quad\hbox{horizontally},\qquad
 a^{-2}q^{1/m}(m-1)q
 \quad\hbox{vertically}.
\]
Consequently,
\begin{equation}\label{eq:link-margin}
 \Ric_{G(a,q,J)}>(m-1)G(a,q,J)
 \quad\Longleftrightarrow\quad
 a^2<q^{1/m}
       \min\left\{1-\frac{2(q-1)}{m-1},\,q\right\}.
\end{equation}
Every \(V_l\) is preserved by the positive Laplacian of \(G(a,q,J)\),
whose restriction is
\begin{equation}\label{eq:angular-operator}
 A_l(a,q,J)
 =a^{-2}q^{1/m}
       \left[B_l\Id+(q^{-1}-1)(-D_J^2|_{V_l})\right].
\end{equation}
\end{lemma}

\begin{proof}
We suppress \(J\) from the notation.
Relative to \(g_0\), the metric \(\widetilde g_q\) has eigenvalues \(1\)
on the \((m-1)\)-dimensional horizontal space and \(q\) on the Hopf
direction.  Its volume form is \(\sqrt q\,d\operatorname{vol}_{g_0}\).
Multiplication of the metric by \(a^2q^{-1/m}\) multiplies this volume
form by \(a^m q^{-1/2}\), proving the volume assertion.

Here is a direct curvature calculation.
Let \(\nabla\) be the round connection and set
\(\phi X=\nabla_X\xi\).  Tangent projection of the ambient derivative
gives
\[
 \phi X=JX+\eta(X)x,\qquad
 \phi^*=-\phi,\qquad \phi\xi=0,\qquad
 \phi^2=-\Id+\xi\otimes\eta,
\]
and differentiation gives
\[
 (\nabla_X\eta)(Y)=g_0(\phi X,Y),\qquad
 (\nabla_X\phi)Y=\eta(Y)X-g_0(X,Y)\xi.
\]
Put \(\delta=q-1\), which is constant on the sphere.
If \(C=\widetilde\nabla-\nabla\) is the difference between the
connections of \(\widetilde g_q\) and \(g_0\), the Koszul formula yields
\begin{align*}
 2\widetilde g_q(C(X,Y),Z)
 &=(\nabla_X\widetilde g_q)(Y,Z)
   +(\nabla_Y\widetilde g_q)(X,Z)
   -(\nabla_Z\widetilde g_q)(X,Y)\\
 &=2\delta\left[\eta(X)g_0(\phi Y,Z)
                    +\eta(Y)g_0(\phi X,Z)\right].
\end{align*}
The vectors \(\phi X,\phi Y\) are horizontal, so
\[
 C(X,Y)=\delta\bigl(\eta(X)\phi Y+\eta(Y)\phi X\bigr).
\]
In particular, \(\tr[X\mapsto C(X,Z)]=0\).
Tracing the connection change formula for curvature therefore gives
\[
 (\Ric_{\widetilde g_q}-\Ric_{g_0})(Y,Z)
 =\tr\bigl[X\mapsto(\nabla_XC)(Y,Z)\bigr]
  -\tr\bigl[X\mapsto C(Y,C(X,Z))\bigr].
\]
These are endomorphism traces and can be computed in a round
orthonormal frame.
Using the identities for \(\phi\), the four terms obtained by
differentiating \(C\) contribute, in order,
\[
 -\delta(g_0-\eta^2),\quad
 \delta(m-1)\eta^2,\quad
 -\delta(g_0-\eta^2),\quad
 \delta(m-1)\eta^2.
\]
Thus the first trace is
\(\delta(-2g_0+2m\eta^2)(Y,Z)\).
For the second trace, one has
\[
 C(Y,C(X,Z))
 =\delta^2\eta(Y)
   \bigl[-\eta(X)Z-\eta(Z)X+2\eta(X)\eta(Z)\xi\bigr],
\]
whose trace in \(X\) is
\(-(m-1)\delta^2\eta(Y)\eta(Z)\).
Since \(\Ric_{g_0}=(m-1)g_0\), this proves the tensor identity
\[
 \Ric_{\widetilde g_q}
 =\bigl((m-1)-2(q-1)\bigr)(g_0-\eta^2)
   +(m-1)q^2\eta^2.
\]
The squared length of \(\xi\) in \(\widetilde g_q\) is \(q\).
Hence the vertical eigenvalue of its Ricci endomorphism is
\((m-1)q\), while the horizontal eigenvalue is
\((m-1)-2(q-1)\).
Constant metric scaling leaves the covariant Ricci tensor unchanged
and divides its endomorphism eigenvalues by the scale.
This proves the stated eigenvalues and \eqref{eq:link-margin}.

Finally, the inverse metric is
\[
 G(a,q,J)^{-1}
 =a^{-2}q^{1/m}
     \left[g_0^{-1}+(q^{-1}-1)\xi\otimes\xi\right].
\]
Because its volume form is a spatially constant multiple of round
volume, its divergence on vector fields is the round divergence.
The field \(\xi\) has zero round divergence, so
\(\operatorname{div}_{g_0}((D_Ju)\xi)=D_J^2u\).
It follows that
\[
 \Delta_{G(a,q,J)}^{+}u
 =a^{-2}q^{1/m}
       \left[\Delta_{g_0}^{+}u+(q^{-1}-1)(-D_J^2u)\right].
\]
Restricting this identity to \(V_l\) gives
\eqref{eq:angular-operator}.
\end{proof}
